\section{Results}
\label{sec:results}

\subsection{ASL: baselines and sequence models}
\label{sec:results-asl}

\begin{table}[t]
  \centering
  \caption{ASL results (\numAslClasses{} signs). Validation and test sets contain 3 and 4 signers
  absent from training. Sequence models: mean $\pm$ standard deviation over three seeds, 50
  epochs. The last column is the standard deviation of top-1 accuracy across the four test
  signers. All values in \%.}
  \label{tab:main}
  \small
  \begin{tabular}{lcccccc}
\toprule
 & & \multicolumn{1}{c}{Validation} & \multicolumn{4}{c}{Test (4 unseen signers)} \\
\cmidrule(lr){3-3}\cmidrule(lr){4-7}
Model & Params & Top-1 & Top-1 & Top-5 & Macro-F1 & Signer s.d. \\
\midrule
Chance & -- & 0.40 & 0.40 & 2.0 & -- & -- \\
Statistics + log.\ regression & -- & 47.9 & 42.1 & 69.4 & 41.9 & -- \\
Statistics + LightGBM & -- & 45.2 & 33.8 & 61.5 & 33.4 & -- \\
\midrule
BiGRU & 1.43\,M & 65.8 $\pm$ 0.5 & 63.0 $\pm$ 0.5 & 84.0 $\pm$ 0.3 & 62.3 $\pm$ 0.6 & 3.9 $\pm$ 0.4 \\
Conv-Transformer & 2.16\,M & 68.2 $\pm$ 0.1 & 64.8 $\pm$ 0.4 & 84.0 $\pm$ 0.3 & 64.1 $\pm$ 0.4 & 3.3 $\pm$ 0.2 \\
\bottomrule
\end{tabular}

\end{table}

\Cref{tab:main} reports the main comparison. Temporal statistics are already informative: a
logistic regression reaches \numAslLogregTest{}\,\% top-1 on unseen signers, about
\numAslLogregChance{} times chance level. Gradient-boosted trees fit the training signers better
but generalise worse (\numAslLgbmTest{}\,\% on test, \numAslLgbmDrop{} points below their
validation score, against \numAslLogregDrop{} points for the linear model). With only 14
training signers, flexible models on hand-crafted statistics capture signer-specific cues.

